% ── Databázový model ──────────────────────────────────────────
\chapter{Databázový model}

        Databázový model tvorí jadro aplikácie a~určuje štruktúru dát, väzby medzi
entitami aj pravidlá integrity. Implementovali sme ho v~SQLAlchemy ORM a~namapovali
na tabuľky SQLite. Pri návrhu sme vychádzali z~potreby prepojiť projektové riadenie, správu
dokumentov a~AI~komunikáciu v~jednom konzistentnom priestore bez zbytočnej duplicity.

        Centrálna je entita projektu, na ktorú sa viažu procesné dáta (úlohy,
plánovanie, notifikácie), obsahové dáta (dokumenty) a~komunikačné dáta (chatové
relácie a~správy). Model kombinuje väzby jeden-ku-viacerým aj viacerí-ku-viacerým.
Hierarchické vzťahy sme riešili priamo cudzími kľúčmi, zatiaľ čo flexibilné prepojenia
(napr. štítky alebo závislosti úloh) sme realizovali cez asociačné tabuľky.

        Životný cyklus údajov sme riadili politikami \texttt{CASCADE} a~\texttt{SET NULL}.
Pri zmazaní nadradeného záznamu sa závislé dáta odstránia tam, kde nemajú samostatný
význam; v~prípadoch, keď je potrebné zachovať históriu, sa referencia nastaví na
\texttt{NULL}. Tým sme znížili riziko osirelých záznamov a~zachovali auditovateľnosť.

        Integritu sme doplnili obmedzeniami \texttt{NOT NULL}, \texttt{UNIQUE} a~kompozitnými
unikátnymi kľúčmi, ktoré bránia duplicitným väzbám v~asociačných tabuľkách. Väčšina
atribútov je relačná (identifikátory, referencie, stavy, časové značky), pričom JSON
polia sme použili len tam, kde bola potrebná vyššia flexibilita, napríklad pre metadáta,
práva alebo citácie odpovedí AI.

        Celkovo sme model navrhli ako rozšíriteľný: jadrové tabuľky, asociačné
tabuľky a~konfiguračné entity umožňujú pridávať nové funkcie inkrementálne bez
zásadných zásahov do existujúcej schémy.
